\documentclass[12pt,a4paper]{article}

% --- Настройка языка и шрифтов для LuaLaTeX ---
\usepackage{fontspec}
\defaultfontfeatures{Ligatures=TeX,Scale=MatchLowercase}

\setmainfont{Alice-Regular}[
    Path=font/,
    Extension=.ttf,
    BoldFont = Alice-Regular,
    ItalicFont = Alice-Regular,
    BoldItalicFont = Alice-Regular,
    BoldFeatures = {FakeBold=1.6},
    ItalicFeatures = {FakeSlant=0.2},
    BoldItalicFeatures = {FakeBold=1.6,FakeSlant=0.2}
]

\newfontfamily\cyrillicfont{Alice-Regular}[
    Path=font/,
    Extension=.ttf,
    BoldFont = Alice-Regular,
    ItalicFont = Alice-Regular,
    BoldItalicFont = Alice-Regular,
    BoldFeatures = {FakeBold=1.6},
    ItalicFeatures = {FakeSlant=0.2},
    BoldItalicFeatures = {FakeBold=1.6,FakeSlant=0.2}
]

\usepackage{polyglossia}
\setmainlanguage{russian}
\setsansfont{TeX Gyre Heros}
\setmonofont{TeX Gyre Cursor}
\AtBeginDocument{\renewcommand{\textnumero}{{\sffamily №}}}

\usepackage{graphicx}
\usepackage{amsmath, amsfonts, amssymb, amsthm}
\usepackage{mathtools}
\usepackage[protrusion=false]{microtype}
\usepackage{unicode-math}
\usepackage{array}
\usepackage{float}
\usepackage{caption}
\usepackage{enumitem}

\usepackage[paper=letterpaper,
            top=20mm,
            bottom=20mm,
            left=10mm,
            right=10mm,
            footskip=0.4in 
]{geometry}

% --- Колонтитулы ---
\usepackage{fancyhdr}
\pagestyle{fancy}
\fancyhf{}
\setlength{\headheight}{15pt}

\renewcommand{\headrulewidth}{0.4pt}
\renewcommand{\footrulewidth}{0.4pt}

\newcommand{\DocNumber}{1}
\newcommand{\DocTitle}{Перспективы развития микропроцессорной техники}
\newcommand{\Discipline}{ЭВМ и периферийные устройства}

\title{
	Министерство науки и высшего образования 
    
    Российской Федерации
    
    Национальный исследовательский университет «МИЭТ»
    
    Институт интегральной электроники имени академика 
    
    К. А. Валиева

	\vspace{2em}
}
\author{
}
\date{}

\begin{document}

\maketitle
\thispagestyle{fancy}

\begin{center}
	\Large Доклад \textnumero~$\DocNumber$\\
	по дисциплине «\Discipline»\\
	«\DocTitle»
\end{center}

\vspace{15em}

\begin{flushright}
	Выполнил студент группы ЭН-$35$\\
	Андреев И. А.
\end{flushright}

\cfoot{Москва 2026}

\pagebreak

\fancyhead[L]{\DocTitle}
\fancyhead[R]{\Discipline}
\fancyfoot[C]{\thepage}

\begin{center}
    \Large\textbf{Доклад \textnumero~$\DocNumber$}\\
    \large\textbf{\DocTitle}
\end{center}

\section{Состояние отрасли в данный момент и объемы производства}

Текущее состояние производства микроэлектроники в России 
характеризуется глубочайшей структурной асимметрией между 
возможностями схемотехнического проектирования по концепции 
fabless и возможностями физической фабрикации кристаллов на 
кремнии в рамках модели foundry.

Исторически российские дизайн-центры, такие как МЦСТ с архитектурой 
«Эльбрус» и «Байкал Электроникс», успешно создавали микропроцессорные 
ядра по проектным нормам 
$28$, $16$ и $12$ нанометров. Однако в рамках глобального разделения 
труда физическое производство полностью размещалось на мощностях 
тайваньской фабрики TSMC. Собственной передовой фабрикации в стране 
создано не было.

Если обратиться к реальным производственным площадкам, действующим 
на территории Российской Федерации сегодня, то их возможности 
ограничены зрелыми технологическими нормами. Флагманом отечественной 
микроэлектроники остается зеленоградское ПАО «Микрон», располагающее 
линиями на кремниевых пластинах диаметром $200$ миллиметров с 
топологическими нормами $180$ и $90$ нанометров, причем устойчивым 
коммерческим процессом является узел $130$--$90$ нанометров. Выпускаемая 
номенклатура функционально ограничена: это чипы для платежных карт 
«Мир», транспортных карт, RFID-метки, дискретные полупроводники и 
микроконтроллеры общего назначения. Предприятие испытывает предельную 
загрузку мощностей, дефицит импортных фоторезистов высокой степени 
очистки и физический износ литографических сканеров глубокого 
ультрафиолета. Фактический коэффициент полезного использования 
оборудования оценивается специалистами на уровне $30$\% из-за 
сложностей с сервисным обслуживанием и технологическими газами.

Второй зеленоградской площадкой является ООО «НМ-Тех», развернутое 
на базе имущественного комплекса завода «Ангстрем-Т». Предприятие 
находится на этапе пусконаладки и постепенного масштабирования 
линий на нормы $130$ и $90$ нанометров на пластинах $200$ миллиметров, 
выпуская драйверы питания, микроконтроллеры и интерфейсные схемы. 
Главной уязвимостью здесь выступает опора на подержанное западное 
оборудование, закупленное до $2022$ года, и острая нехватка оригинальных 
запасных частей. В свою очередь, ОАО «Интеграл» работает в 
диапазоне норм $350$--$180$ нанометров, специализируясь на силовой 
электронике, интерфейсных БИС и радиационно-стойкой ЭКБ для 
специальной техники. Выпуск современных микропроцессоров на этих 
мощностях физически невозможен.

Отвечая на принципиальный вопрос о достаточности объемов производства, 
следует признать, что продукция выпускается в критически недостаточном 
количестве. Суммарный годовой объем выпуска отечественных 
микроконтроллеров и интегральных схем всех типов внутри страны 
составляет порядка $7,5$--$8,5$ миллионов штук. В то же время годовая 
потребность одной лишь российской гражданской и промышленной 
экономики исчисляется десятками миллионов микросхем для автопрома, 
приборостроения, интеллектуальных приборов учета и промышленной 
автоматизации.

Что касается высокопроизводительных микропроцессоров серверного и 
десктопного уровня, потребность в которых составляет сотни тысяч 
штук в год, то их физическое производство на территории России в 
настоящее время равно абсолютному нулю. Доля отечественных 
производителей на внутреннем рынке полупроводников составляет 
лишь около $28$\% в денежном выражении, а в сегменте микропроцессоров 
высокой степени интеграции не превышает одного-двух процентов. 
Физический предел в $90$ нанометров делает невозможным локальное 
изготовление процессоров для серверов, персональных компьютеров и 
телекоммуникационного оборудования.

\section{Воздействие санкций на технологический цикл}

Введенные после $2022$ года масштабные международные санкции нанесли 
системный деструктивный удар по отечественной микроэлектронике сразу 
по трем критическим направлениям технологического цикла.

Первым ударом стала тотальная фаундри-блокада. Отказ тайваньской 
фабрики TSMC от выполнения контрактов привел к физической остановке 
серийного выпуска передовых российских микропроцессоров линеек 
«Эльбрус-8СВ», «Эльбрус-16С», «Байкал-М» и «Байкал-S». Сотни тысяч 
уже изготовленных и оплаченных кремниевых кристаллов оказались 
заблокированы на зарубежных складах без возможности отгрузки. 
Попытки перенести заказы на альтернативные азиатские фабрики, 
прежде всего на китайскую SMIC, натолкнулись на жесткий барьер 
экстерриториального правила США о прямом иностранном продукте. 
Фабрика SMIC критически зависит от американского оборудования 
Applied Materials, Lam Research и KLA, и под угрозой попадания 
в блокирующие санкционные списки SDN отказывается принимать 
российские заказы на проектные нормы тоньше $40$ нанометров.

Вторым ударом стала изоляция отечественных дизайн-центров от мировых 
систем автоматизированного проектирования. Лидеры глобального рынка 
САПР --- Synopsys, Cadence и Siemens EDA --- отозвали лицензии и 
заблокировали доступ к обновлениям и библиотекам стандартных ячеек. 
Без этих программных комплексов невозможно выполнить физический 
дизайн СБИС, синтез многомиллионных межсоединений, статический 
временной анализ и топологическую верификацию. Существующие 
российские программные продукты закрывают лишь локальные задачи 
проектирования печатных плат и простейших ПЛИС, не поддерживая 
нормы глубокого субмикрона.

Третьим критическим направлением удара стало эмбарго на поставку 
технологического оборудования и высокочистых материалов. 
Российские чистые комнаты были полностью отрезаны от оборудования 
для литографии глубокого и экстремального ультрафиолета компаний 
ASML и Nikon. Были заблокированы легальные поставки кремниевых 
пластин диаметром $300$ миллиметров, полировальных суспензий, 
фотошаблонов и сверхчистых фоторезистов от мировых монополистов 
Shin-Etsu, Tokyo Ohka Kogyo и Merck. В результате действующие 
отечественные полупроводниковые фабрики оказались в условиях 
тяжелого ресурсного и технологического дефицита.

\section{Пути преодоления санкций и оценка их успешности}

В ответ на внешние ограничения российское государство и промышленность 
развернули комплекс мер структурной адаптации, сосредоточенный на 
трех ключевых механизмах: логистической реконфигурации, переходе на 
открытые архитектуры и развитии собственного электронного машиностроения.

Первый механизм выразился в выстраивании каналов параллельного импорта 
через юрисдикции Китая, Гонконга, Турции и Объединенных Арабских 
Эмиратов. Это позволило частично компенсировать первоначальный шок и 
восстановить ввоз широкой номенклатуры микросхем западных 
производителей. Однако параллельный импорт породил комплекс 
тяжелейших проблем. Закупочные цены на компоненты выросли в $2$--$5$ раз 
по сравнению с базовыми. Введение вторичных санкций привело к 
блокировкам трансграничных платежей и затягиванию поставок на месяцы. 
Кроме того, в партиях лавинообразно возросла доля бракованной, 
восстановленной и перемаркированной продукции, что несет критические 
риски отказов аппаратуры.

Вторым направлением стал переход российских разработчиков на открытую 
международную архитектуру процессорных команд RISC-V. Утрата доступа к 
лицензированию коммерческих ядер британской архитектуры ARM вынудила 
такие дизайн-центры, как Yadro, Syntacore и CloudBEAR, форсировать 
разработку собственных вычислительных ядер на базе RISC-V. Это 
стратегически верный шаг, полностью снимающий патентные и лицензионные 
риски на стадии проектирования.

Третье направление заключается в запуске масштабных государственных 
программ развития электронного машиностроения и химических материалов 
с объемом финансирования свыше ста миллиардов рублей, что не является 
серьёзной суммой для данной сферы. В рамках этих 
программ институтами РАН и промышленными партнерами собран первый 
экспериментальный макет литографа на топологическую норму $350$ 
нанометров, развернуты опытно-конструкторские работы по созданию 
промышленного степпера на $130$ нанометров со сроком демонстрации к 
$2027$ году, ведутся исследования рентгеновской литографии, а также 
освоено производство ряда особо чистых технологических газов.

Оценивая успешность преодоления санкций, ее следует охарактеризовать 
как частично успешную в сегменте зрелой микроэлектроники и крайне 
тяжелую в сфере вычислительных микропроцессоров. Полупроводниковое 
производство требует непрерывной цепочки из десятков сложнейших 
типов оборудования --- установок плазмохимического травления, ионной 
имплантации, планаризации и эпитаксии. Воссоздать всю эту глобальную 
индустрию силами одной страны за несколько лет невозможно.

\section{Прогноз состояния микроэлектроники по временным горизонтам}

Анализ производственных циклов и сроков реализации 
научно-исследовательских работ позволяет сформировать научно 
обоснованный прогноз динамики отрасли по временным горизонтам.

В краткосрочной перспективе, на горизонте полугода, фабрики «Микрон» и 
«НМ-Тех» продолжат функционировать на технологических нормах $180$--$90$ 
нанометров, расходуя имеющиеся складские запасы импортных расходных 
материалов и комплектующих. В сегменте процессоров для персональных 
компьютеров и серверов сохранится острейший дефицит, а гражданский 
сектор будет практически полностью зависеть от серых каналов импорта 
на фоне постоянного ужесточения банковского комплаенс-контроля.

Через год ожидается завершение основных пусконаладочных работ на заводе 
«НМ-Тех» по топологиям $130$ и $90$ нанометров. Это позволит увеличить 
совокупный годовой выпуск отечественных микроконтроллеров и смарт-карт 
до $10$--$12$ миллионов штук, а также внедрить первые партии локализованных 
химических реактивов. Однако отечественные вычислительные процессоры в 
серийном производстве по-прежнему будут отсутствовать, а ключевым 
вызовом станет нарастающий износ импортного парка литографических 
сканеров.

На горизонте двух лет планируется ввод в опытную эксплуатацию первых 
отечественных литографических установок с разрешением $350$ нанометров и 
начало стендовых испытаний оптического тракта литографа на $130$ 
нанометров. В этот период появятся первые промышленные серии 
отечественных микроконтроллеров на архитектуре RISC-V, физически 
произведенные на зеленоградских фабриках, хотя трудности 
масштабирования сложных оптических узлов и эксимерных лазеров останутся 
крайне острыми.

В пятилетней перспективе, к $2029$--$2031$ годам, отрасль попытается 
освоить проектную норму $65$ нанометров на пластинах $200$ миллиметров и 
развернуть опытное оборудование для пластин диаметром $300$ миллиметров. 
В этот период станет возможным заметное импортозамещение 
микроконтроллеров в энергетике, приборостроении и автопроме. При этом 
технологическое отставание от мирового фронтира, 
осваивающего нормы $2$ и $1,4$ нанометра, законсервируется на уровне 
$15$--$20$ лет.

За пределами пятилетнего горизонта, после $2031$ года, прогнозируется 
появление первых опытных отечественных установок суб-$28$-нанометрового 
диапазона, включая рентгеновские литографы, и формирование замкнутого 
суверенного контура микроэлектроники для критической информационной 
инфраструктуры. Главным стратегическим вызовом этого периода станет 
окончательный переход от многовекторной зависимости от Запада к жесткой 
технологической зависимости от китайских производителей оборудования и 
контрактных фабрик.

\section{Влияние состояния отрасли на экономику страны в целом}

Влияние состояния микроэлектроники на экономику России квалифицируется 
как существенное, глубокое и долгосрочное. Дефицит полупроводников 
напрямую трансформирует структуру ключевых отраслей-потребителей.

В автомобильной промышленности нехватка микроконтроллеров для блоков 
управления двигателем, систем ABS и курсовой устойчивости уже привела к 
временному упрощению регламентов безопасности. Массовый переход на 
безальтернативную китайскую электронную компонентную базу 
сопровождается ростом себестоимости электронных блоков на $30$--$60$\% и 
потерей схемотехнической самостоятельности отечественного автопрома.

В секторе телекоммуникаций и центров обработки данных недоступность 
быстродействующих ПЛИС, цифровых сигнальных процессоров и серверных 
CPU привела к фактической остановке проектов развертывания сетей связи 
$5G$ в России. Серийное производство отечественных базовых станций 
сдвинуто за пределы $2030$ года. Корпоративный сектор столкнулся с 
каннибализацией оборудования и лавинообразным ростом затрат на 
вычислительную инфраструктуру, что тормозит развитие технологий 
искусственного интеллекта и облачных сервисов.

В оборонно-промышленном комплексе микроэлектроника выступает абсолютным 
приоритетом: сюда направляется до $70$--$80$\% доступной компонентной базы 
специального назначения. Складские запасы и параллельные поставки 
позволяют закрывать первоочередные задачи, однако усложнение логистики 
и риски контрафакта вызывают задержки при серийном изготовлении 
высокоточных систем и бортовой радиолокации.

В сфере промышленной автоматизации и топливно-энергетического комплекса 
импортозамещение контроллеров АСУ ТП сдерживается отсутствием 
отечественных прецизионных АЦП и скоростных интерфейсных микросхем. 
Предприятия вынуждены замещать западное оборудование решениями 
производителей из КНР, попадая в прямую зависимость от китайских 
сервисных платформ и закрытых протоколов.

\section{Влияние перспектив отрасли на личные профессиональные планы}

Для меня как студента Национального исследовательского университета 
«МИЭТ» института ИнЭл, обучающегося по инженерному профилю 
микроэлектроники, происходящие процессы носят прямо определяющий 
характер. Если ранее стандартная карьерная траектория предполагала 
разработку блоков для фабрикации на зарубежных заводах под нормы $28$ 
или $16$ нанометров, то сегодня мой профессиональный фокус смещается по 
трем конкретным направлениям.

Первое направление --- это схемотехническое проектирование на открытой 
архитектуре RISC-V и концепция архитектурной компенсации ограничений 
кремния. В условиях, когда отечественное физическое производство будет 
ограничено нормами $90$, $65$ и в лучшем случае $28$ нанометров, инженерная 
задача состоит в том, чтобы добиваться роста производительности не 
уменьшением транзистора, а развитием архитектур, упущенных западными
компаниями. Также возможно применение китайского опыта реверс-инжиниринга
опытных образцов процессоров и литографов.

Вторым направлением становится специализация на входном контроле. 
Массовый приток электронных 
компонентов по параллельному импорту создал огромный спрос на 
специалистов, владеющих методами рентгеновской микротомографии, 
сканирующей акустической микроскопии, анализа вскрытых кристаллов и 
верификации топологий. Умение безошибочно выявлять контрафакт, следы 
деградации или недокументированные аппаратные закладки в 
микропроцессорах является сегодня одной из самых критически 
востребованных и высокооплачиваемых инженерных компетенций.

Третье направление связано с участием в научно-исследовательских и 
опытно-конструкторских работах по отечественному оборудованию. 
Будущее микроэлектроники сегодня решается на стыке физики твердого 
тела, плазмохимии, лазерной оптики и вакуумной техники. Разработка 
отечественных источников излучения, прецизионных столов литографов, 
химических реакторов и технологий сверхчистых материалов в лабораториях 
МИЭТ и зеленоградских предприятиях --- это та сфера, где 
инженер может внести реальный вклад в технологический суверенитет 
страны.

\section{Заключение}

Подводя итог докладу, хочу подчеркнуть ключевую мысль: иллюзия того, 
что микроэлектронику можно быстро восстановить за счет одного лишь 
государственного финансирования или полностью заменить параллельным 
импортом, развеяна реальностью. Отечественная полупроводниковая отрасль 
вступила в эпоху долгой и глубокой структурной перестройки. Наш 
единственный реалистичный путь --- создавать отечественные средства 
производства с нуля или искать договорённости с партнёрами в иных 
плоскостях.

\end{document}
